% from ega4-2-fr.tex
\oldpage[IV-2]{186}
que $A$ est \emph{intègre} et est un sous-anneau de $B$. On sait (1.3.11)
que $B$
est quotient d'un anneau local de la forme $C_{\mathfrak p}$, où
$C=A[T_1,\ldots,T_n]$ est une algèbre de polynômes et $\mathfrak p$ un
idéal premier de $C$ au-dessus de l'idéal maximal $\mathfrak m$ de $A$. En
vertu de (7.1.7), il suffit de prouver que $C_{\mathfrak p}$ est
formellement équidimensionnel, autrement dit on peut se borner au cas où
$B=C_{\mathfrak p}$.

Posons $A'=\widehat A$, et
$C'=A'[T_1,\ldots,T_n]=C\otimes_AA'$; il y a un idéal premier unique
$\mathfrak p'$ de $C'$ au-dessus de l'idéal maximal $\mathfrak mA'$ de
$A'$, donc au-dessus de $\mathfrak p$; posons $B'=C'_{\mathfrak p'}$;
l'homomorphisme $B\to B'$ est local et fait de $B'$ un $B$-module plat. On
sait alors que $\widehat {B'}$ est un $\widehat B$-module plat (Bourbaki,
\emph{Alg. comm.}, chap. III, § 5, n\textsuperscript{o} 4, prop. 4); comme
$\widehat {B'}$ est caténaire (5.6.4), il suffira de prouver que
$\widehat {B'}$ est \emph{équidimensionnel} pour en déduire que $\widehat B$ l'est
aussi (7.1.3, (i)), ce qui achèvera la démonstration.

Or, $A'$ est quotient d'un anneau régulier par le théorème de Cohen (0,
19.8.8), donc il en est de même de $B'$ (0, 17.3.9); en vertu de (7.1.5),
pour montrer que $B'$ est formellement équidimensionnel, il suffit de
prouver que $B'$ est équidimensionnel; et pour cela, il suffit de montrer
que $C'$ est équidimensionnel, puisque $C'$ est quotient d'anneau régulier,
donc caténaire (5.6.4). Or les idéaux premiers minimaux de $C'$ sont les
idéaux $\mathfrak q'C'$, où $\mathfrak q'$ parcourt les idéaux premiers
minimaux de $A'$ (5.5.3), et on a
$C'/\mathfrak q'C'=(A'/\mathfrak q')[T_1,\ldots,T_n]$. Comme $A'$ est
équidimensionnel par hypothèse (puisque $A$ est intègre), il en est donc de
même de $C'$ par (5.5.4). C.Q.F.D.
\end{proof}

\begin{definition}[7.1.9]
\label{IV.2.s7.1.9-fr}
Lorsque les conditions équivalentes de (7.1.8) sont vérifiées, on dit que
$A$ est un anneau \emph{formellement caténaire}.
\end{definition}

Pour un anneau local noethérien, il revient donc au même de dire qu'il est
formellement équidimensionnel ou formellement caténaire et
équidimensionnel, en vertu de (7.1.2).

\begin{proposition}[7.1.10]
\label{IV.2.s7.1.10-fr}
Tout anneau local $A$ quotient d'un anneau local de Cohen--Macaulay est
formellement caténaire.
\end{proposition}

\begin{proof}
Cela résulte aussitôt de (7.1.8) et de (7.1.5) appliqué aux anneaux
quotients intègres de $A$.
\end{proof}

\begin{proposition}[7.1.11]
\label{IV.2.s7.1.11-fr}
\begin{enumerate}
\item[(i)] Un anneau local noethérien formellement caténaire est universellement
caténaire.
\item[(ii)] Si $A$ est un anneau local noethérien formellement caténaire, tout
anneau local $B$ qui est une $A$-algèbre essentiellement de type fini est
formellement caténaire.
\end{enumerate}
\end{proposition}

\begin{proof}
L'assertion (ii) résulte aussitôt de la condition c) de (7.1.8) et du fait
que si un anneau local $C$ est une $B$-algèbre essentiellement de type
fini, $C$ est aussi une $A$-algèbre essentiellement de type fini (1.3.9).
Pour prouver (i), notons d'abord qu'en vertu de la condition b) de (7.1.8),
un anneau local noethérien formellement caténaire $A$ est caténaire,
puisque les quotients de $A$ par ses idéaux premiers minimaux le sont
(7.1.4). Si maintenant $E$ est une $A$-algèbre de type fini, il résulte de
ce qui précède et de (ii) que tout anneau local de $E$ en un idéal premier
est caténaire, donc que $E$ est caténaire. Donc $A$ est universellement
caténaire.
\end{proof}

